We organize a collection of fixed-shift prime-pair investigations into seven modules: arithmetic foundations, analytic cancellation, fixed atoms, actual-occurrence structures, nonautonomous dynamics, the twin-prime endpoint, and auxiliary geometric lifts. The purpose is to distinguish established structures from transfers that still require arithmetic input, while keeping the objects, quantifiers, and normalizations consistent. The analytic core is an arithmetic operator retaining unit masks, a smooth kernel, and diagonal deletion. We restate and prove its exact positive-form decomposition, norm and row-energy bounds, and the restrictions that the actual coefficient norms impose on coefficient-blind methods. At interval scale $x$ and prime-modulus scale $Q=x^{1/3}$, a low-frequency compression yields $(1/4+o(1))Q^2$ negative directions. An exact Kloosterman representation explains why one-sided short support does not automatically supply a two-sided critical-range estimate. For an explicitly constructed smooth partition, the collective boundary contribution is bounded by $xQ^2x^{-11/64+o(1)}$. On the dynamical side, the Chinese remainder theorem and Abel summation give a linear variance bound and Haar-almost-everywhere infinite recurrence for the actual moving sieve events, while leaving the return to the distinguished arithmetic seed $0$ unproved. These results establish neither the required signed interior cancellation nor the twin-prime conjecture. They provide a source-traceable synthesis of structures and obstructions that separates two unresolved bridging problems from the available results.
